\section{Related Work}
\label{sec:related-work}

\subsection{Model Side Channels and Compiler-Centered ML Security}

Prior attacks establish that a model's execution can reveal much more than its final prediction.
CSI NN recovers neural-network architecture parameters from electromagnetic traces~\cite{batina2019csinn}; Cache Telepathy infers DNN architectures through shared-cache behavior~\cite{yan2020cachetelepathy}; and DeepSniffer reconstructs architectural hints from accelerator execution signals~\cite{hu2020deepsniffer}.
These systems demonstrate the exploitability of model-dependent timing, memory, and microarchitectural signatures.
Our objective is complementary: we do not construct a complete model-extraction attack.
We ask which signatures cross a compiler boundary, which are created or removed there, and whether an individual lowering specification preserves a declared observation policy.

Compilers can also threaten ML systems through integrity failures.
Chen et al. show that inconsistencies between eager and compiled execution can be weaponized as model backdoors~\cite{chen2026compilerbackdoor}.
We use a related eager-versus-compiled comparison, but for confidentiality rather than output integrity: identical public behavior does not suffice when timing, control flow, or addresses remain distinguishable.
Together, these lines of work motivate treating the compiler as part of the security boundary rather than as a semantics-neutral implementation detail.

\subsection{Finding Constant-Time Violations}

Constant-time analysis spans source, intermediate, and binary levels. ct-verif reduces constant-time verification to safety properties over a self-composed program~\cite{almeida2016ctverif}.
Binsec/Rel performs relational symbolic execution directly on binaries and demonstrates that compiler choices can create constant-time violations~\cite{daniel2020binsecrel}.
Simon et al. likewise show how mainstream C compilation can invalidate source-level control over security-relevant side effects~\cite{simon2018whatyouc}.
These results preclude a claim that compiler-induced leakage or differential compiler analysis is new in itself.

Our distinction is the workflow and workload boundary.
We combine count-based and taint-based observations on ML kernels, execute eager and classify a channel as oblivious, distinguishable, compiler introduced, or compiler removed at a pinned configuration point.
That four-way attribution is empirical rather than a constant-time proof.
The FCVD-CT layer then supplies the complementary universal check for the inputs and observations represented by its symbolic model.
In particular, it treats control flow, addresses, variable-latency operations, and selected resource counts as explicit semantic observations rather than reducing all leakage to wall-clock time.

\subsection{Security-Preserving Compilation}

The closest formal work studies whether compilation preserves information-flow or noninterference hyperproperties.
Besson et al. formulate information-flow-preserving optimizations and apply translation validation to compiler transformations~\cite{besson2019ifpreservation}.
Namjoshi and Tabajara generalize this direction with witness relations for validating secure compilation~\cite{namjoshi2020witnessing}.
Barthe et al. machine-check a constant-time-preserving C compiler based on CompCert ~\cite{barthe2020ctcompiler}.
More recently, SNIP checks preservation of speculative noninterference and shows that transformations such as register allocation must account for observations introduced by speculation ~\cite{vanderwall2025snip}.

Those systems provide general proof principles or stronger guarantees for a fixed compiler chain.
We make a different engineering tradeoff.
FCVD-CT instantiates leakage-aware translation validation for extensible MLIR dialect semantics and checks a pair of pass specifications with two independent queries: functional refinement and channel-relative noninterference.
This does not amount to a whole-compiler proof; the correspondence between a handwritten specification and its production pass remains trusted until actual before/after IR is validated.
The benefit is that a dialect author can add semantics and an observation rule locally, while falsifying controls expose whether each of the two obligations is genuinely load-bearing.

\texttt{ct\_select} intrinsic added to Clang and LLVM lowers to an optimization barrier that is selected into branch-free machine code, but it protects only the selections a programmer has marked and only against secret-dependent branching.

\subsection{Translation Validation and MLIR Verification}

Translation validation checks each compilation instance instead of verifying the compiler implementation once and for all.
Alive2 applies bounded refinement checking to LLVM transformations~\cite{lopes2021alive2}.
MLIR-TV brings systematic semantic validation to multiple MLIR dialects and lowering boundaries~\cite{wang2024mlirtv}.
First-Class Verification Dialects (FCVD) package dialect semantics so that verification tools can reuse them ~\cite{fcvd2025}; our checker builds directly on that abstraction.
MLIR's multi-level organization~\cite{lattner2021mlir} and Polygeist's C-to-MLIR pipeline~\cite{moses2021polygeist} provide the substrate and the realistic multi-stage evaluation target.

These validators primarily ask whether the target computes a result permitted by the source.
For confidential compilation, that condition is necessary but not sufficient: a target may compute the same value while exposing a secret-dependent branch or address.
Our extension makes leakage observations part of the executable semantics, self-composes the source and target, and accepts a specification only when both the functional and relational obligations hold.
The contribution is therefore not a new notion of secure compilation, but a concrete, coverage-accounted realization for MLIR lowering specifications coupled to empirical attribution on ML workloads.
